\section{Implementation Interfaces and Lifecycle}
\label{app:interfaces}

This appendix records source-level contracts and operational defaults that are
useful for reproduction but too detailed for the main argument. The lifecycle is
specified directly through role contracts, outcome states, persistence surfaces,
and guard conditions; no additional conceptual diagram is required.

\subsection{Role Contracts}

The Manager front door uses one grounded model decision to emit six structured
axes: \code{CONFIG}, \code{CONTROL}, \code{ROUTE}, \code{LIFETIME},
\code{FAST\_REPLY}, and \code{NAME}. A \code{CONTROL: NO\_DISPATCH} result is
authoritative and fails closed: the bridge may provide a read-only response but may
not enqueue work. The Manager is also the sole writer of the current campaign
stage.

The Planner produces \code{TaskSpec} objects with explicit scope, optional DAG
dependencies, and impact score. It authors and applies stage \code{checklist\_ops};
the Reviewer can send \code{checklist\_feedback} but does not edit the checklist.
The Engineer receives reservation-derived provider fences through
\code{RunnerOptions} and returns a \code{RunnerResult} with a call identifier,
usage-presence flags, cost, and pricing status. The Reviewer returns a
\code{ReviewDecision} whose central fields are listed in
Table~\ref{tab:review-interface}.

\begin{table}[H]
\centering
\footnotesize
\begin{tabularx}{\linewidth}{@{}p{3.3cm} X@{}}
\toprule
\textbf{Reviewer field} & \textbf{Meaning} \\
\midrule
\code{status} & \code{done}, \code{continue}, or \code{blocked}; the mission-level completion verdict \\
\code{progress\_class} & \code{decision}, \code{evidence}, \code{setup\_only}, \code{artifact\_sync\_only}, or \code{none} \\
\code{next\_action} & Concrete action for the next Engineer round \\
\code{operator\_question} & At most one human-facing blocking question \\
\code{scope, checklist} & Structured final-submission gate, evaluated fail-closed \\
\code{skill\_ops, wiki\_ops} & Reviewer-authored retained-state edits \\
\code{backend\_unavailable} & Infrastructure failure, never interpreted as a scientific \code{continue} \\
\bottomrule
\end{tabularx}
\caption{Abridged Reviewer interface from \srcpath{core/models.py} and
\srcpath{reviewer/reviewer\_schema.json}.}
\label{tab:review-interface}
\end{table}

\subsection{Outcome Taxonomy}

Every mission produces an explicit outcome:

\begin{itemize}
  \item \textbf{Terminal:} \code{done}, \code{max\_rounds}, \code{blocked},
    \code{no\_progress}, \code{error}, and \code{budget\_exhausted}.
  \item \textbf{Paused:} \code{paused\_budget},
    \code{paused\_provider\_cooldown}, \code{paused\_provider\_fence},
    \code{infra\_blocked}, \code{research\_incomplete},
    \code{paused\_no\_breakthrough}, and
    \code{exhausted\_current\_methods}.
  \item \textbf{Control transfer:} \code{replan\_requested}; the item returns to
    pending and passes through the normal Planner gate rather than being counted as
    completed or failed.
\end{itemize}

\section{Persistence and Context Management}
\label{app:persistence}

The global state root defaults to \code{\textasciitilde/.argus-skill/} and can be
overridden by \code{ARGUS\_SKILL\_HOME}. Each project is stored under a fingerprinted
subdirectory. Table~\ref{tab:persisted-state} lists the principal files.

\begin{table}[H]
\centering
\footnotesize
\begin{tabularx}{\linewidth}{@{}p{3.0cm} p{2.0cm} X p{1.75cm}@{}}
\toprule
\textbf{Path} & \textbf{Owner} & \textbf{Content} & \textbf{Mutation} \\
\midrule
\code{events.jsonl} & record plane & Canonical typed-event timeline & append-only \\
\code{backlog.jsonl} & scheduler & Pending, running, completed, and superseded tasks & atomic claim / CAS \\
\code{continuous.json} & Manager / scheduler & Standing objective, route, lifetime, and identity & atomic rewrite \\
\code{inbox.jsonl} & operator & Operator messages and answers & append + drain \\
\code{checkpoint.json} & Reviewer & Hard-capped cross-session working state & atomic rewrite \\
\code{daemon.pid} & daemon & Per-project concurrency lock & exclusive \\
\code{daemon.status.json} & daemon & Liveness and handoff status & rewrite \\
\code{skills/} & Reviewer / operator & Versioned reusable skills & versioned \\
\bottomrule
\end{tabularx}
\caption{Principal persisted state. The event tape is the canonical timeline;
other files are working state or materialized views.}
\label{tab:persisted-state}
\end{table}

The checkpoint contains the goal, verified work, tried-and-failed approaches, one
open blocker, the next step, and bounded environment facts. The Engineer proposes a
handoff, and the Reviewer validates and writes the canonical \code{checkpoint.json}.
Python-level item and character caps prevent it from becoming an append-only
transcript. A missing or malformed candidate checkpoint preserves the last valid
state rather than clearing memory.

By default, each role thread is rolled after three rounds on that thread or when the
preceding call reports at least 1.5 million input tokens. The next session is
re-seeded from the checkpoint. These limits are configurable and disable at zero.
The design keeps provider prefix caching useful over short windows while preventing
unbounded compaction loops.

Unattended execution is wrapped by
\srcpath{daemon/life\_worker.py} (roughly 1,800 lines). The worker drains on
\code{SIGTERM} or \code{SIGINT} at a mission boundary. Source upgrades use a
blue/green self-handoff with a bounded generation count and persisted identity
consisting of the objective hash, vertical, and lineage generation. A resume adopts
only a matching identity; otherwise the Manager must perform a fresh division.

\section{Run Records and Provenance}
\label{app:evidence-map}

\subsection{Event Catalog and Call Correlation}

The canonical catalog contains \textbf{112 typed event types} across
\textbf{11 categories}; \textbf{75} event types carry validated payload schemas in
\srcpath{core/event\_catalog.py}. Table~\ref{tab:event-taxonomy} gives the category
counts. The event \code{research.achievement.certified} is the sole authoritative
record of a verified achievement.

\begin{table}[H]
\centering
\footnotesize
\begin{tabularx}{0.86\linewidth}{@{}p{2.3cm} r X@{}}
\toprule
\textbf{Category} & \textbf{Count} & \textbf{Representative content} \\
\midrule
Lifecycle & 40 & Loop, round, mission, and budget-reservation events \\
Skill & 20 & Match, distillation, use, and outcome events \\
Wiki & 17 & Source ingest, page synthesis, promotion, and validation \\
Planner & 9 & Planning start, tasks, waiting contracts, and verdicts \\
Research & 7 & Achievement certification and research status \\
Agent I/O & 5 & Provider-call start, stream, and completion \\
Daemon & 5 & Start, stop, handoff, and status \\
Idea & 3 & Idea lifecycle \\
Provider & 3 & Provider request start, completion, and denial \\
Usage & 2 & Presence-aware usage and cost \\
Operator & 1 & Operator inbox event \\
\bottomrule
\end{tabularx}
\caption{Event taxonomy in the current implementation.}
\label{tab:event-taxonomy}
\end{table}

Every provider call receives a \code{call\_id} at
\srcpath{core/run\_gateway.py}. The identifier is written to the event tape and
returned in the runner result, allowing Reviewer prompts and later audits to locate
the exact execution record. Usage-presence booleans distinguish a genuine zero from
provider-omitted accounting.

\subsection{Credential Handling and Optional Signing}

\srcpath{core/secret\_guard.py} redacts authorization headers, API-key patterns,
provider tokens, and generic key/value secrets before event or artifact copies are
persisted for downstream use. A redaction emits \code{round.secret\_redacted}; a
scan error or uncovered oversized artifact blocks unconditional certification.

An optional Ed25519 signing path in \srcpath{team/result\_provenance.py} supports
isolated-scorer teammate results. It is off by default and is not used for the
results reported here.

\subsection{Artifact-Digest Records}

The nanoGPT speedrun row points to logical project
\nolinkurl{nanogpt-speedrun-h100} and artifact
\nolinkurl{candidate_mlp_post_only_hybrid_n10_20260617T152904Z}. Its verifier
records \code{valid=true}, $N{=}10$, validation loss $3.2776$, one-sided
$p(\mathrm{mean}<3.28)=0.004007$, training time $79.77\pm0.06$ seconds, and
\code{seal=ok}. The nanochat B200 row points to logical project
\nolinkurl{nanochat-mission-b200} and artifact
\nolinkurl{a374_tailweighted_boundary_mixed_readout_data_bundle}; its frozen scorer
exits 0 with \code{MEAN\_VAL\_BPB=0.963634} and a valid SM100 flash-attention kernel.
The result JSON stores source/result hashes and verifier excerpts, but not the
external bytes.

Both public pages were retrieved at \code{2026-07-14T12:29:51Z}. Their complete
raw-HTML SHA-256 values are recorded in \srcpath{evidence/website\_results.json}
and \srcpath{evidence/paper\_inventory.json} rather than typeset inline.

\subsection{Source Map}

\begin{table}[H]
\centering
\footnotesize
\begin{tabularx}{\linewidth}{@{}p{4.1cm} X@{}}
\toprule
\textbf{Subsystem} & \textbf{Primary source path(s)} \\
\midrule
Runtime composition and roles & \code{manager/}, \code{planner/}, \code{engineer/}, \code{reviewer/}, \code{loop.py} \\
Mission lifecycle and daemon & \code{life/supervisor/\_core.py}, \code{daemon/life\_worker.py}, \code{daemon/handoff.py} \\
State and memory & \code{core/paths.py}, \code{apps/\_runtime.py}, \code{engineer/checkpoint.py}, \code{life/memory.py} \\
Execution and budgets & \code{core/run\_gateway.py}, \code{core/pricing.py}, \code{core/usage.py}, \code{core/provider\_quota.py} \\
Reliability guards & \code{engineer/runner.py}, \code{regime\_jump/} \\
Records and integrity & \code{core/event\_catalog.py}, \code{core/secret\_guard.py}, \code{skills/provenance.py} \\
Cockpit and API & \code{webapi/server.py}, \code{frontend/tui/}, \code{frontend/web/} \\
Public results & \code{technical\_report/evidence/website\_results.json}, \code{paper\_inventory.json} \\
\bottomrule
\end{tabularx}
\caption{Repository map for implementation and reported results.}
\label{tab:source-map}
\end{table}

\section{Reliability, Budgets, and Deployment}
\label{app:reliability}

The guards in Table~\ref{tab:guards} define the long-running execution envelope.
They are task-agnostic: they constrain activity and resource use, while the Reviewer
remains responsible for scientific judgment.

\begin{table}[H]
\centering
\footnotesize
\begin{tabularx}{\linewidth}{@{}p{2.75cm} X p{1.3cm} p{1.55cm}@{}}
\toprule
\textbf{Guard} & \textbf{Trigger and action} & \textbf{Default} & \textbf{Live-call interrupt} \\
\midrule
Backend failure & Consecutive failures $\rightarrow$ backoff, then error & 2 / 15 s & no \\
No-progress bail & Consecutive rounds with no Engineer message & 2 & no \\
Decision stall & Consecutive nondecision review classes $\rightarrow$ end & 4 & no \\
Decision-progress budget & Time since last decision/evidence $\rightarrow$ end & 1,800 s & no \\
Soft round escalation & Ask Reviewer to return \code{blocked} & 12 & no \\
Hard round escalation & Force terminal \code{blocked} & 24 & no \\
Effective-progress watchdog & No file change or non-token event & 3,600 s & yes \\
Runner hard-idle & No runner activity & 3,600 s & yes \\
In-round compaction limit & Repeated automatic compactions & 3 & yes \\
\bottomrule
\end{tabularx}
\caption{Default reliability guards in \srcpath{engineer/runner.py}. Thresholds are
environment-overridable and disable at zero where applicable.}
\label{tab:guards}
\end{table}

The canonical \code{resolve\_budget\_caps} path in \srcpath{core/knobs.py} resolves
a per-mission preflight cap of \$30, a per-project daily cap of \$180, a global
daily cap of \$30, and at most two active daemons across projects. A positive global
cap is enforced; setting it to zero disables it. Calls reserve budget before
execution and settle or release the reservation afterward. Unpriced calls and
provider-fence breaches produce explicit blocked events rather than fabricated
costs.

Supervised background jobs maintain their own health cadence and report terminal
outcomes through the operator inbox. An Engineer may explicitly emit
\code{WAIT\_FOR\_SUBAGENT} for a healthy self-watched job; otherwise the request
falls back to a normal reviewed round. The GPU lease tool is explicit rather than
automatic: \code{gpu\_lease run -- <cmd>} releases parked devices, records a lease
for the command lifetime, and restores the keep-alive only when no competing lease
remains.

The intended deployment is a per-project \code{systemd --user} service. It runs the
daemon in the foreground, uses \code{SIGTERM}, allows up to 1,800 seconds for a clean
mission-boundary drain, and restarts under a bounded service policy. Sandboxed
builder roles are confined to the project working directory and may not write the
global state root or provider configuration. The terminal and web cockpits read the
same event tape; mutating commands and sensitive global views require an
authorization token when configured.

\section{Notation and Figure Provenance}
\label{app:notation}

\begin{table}[H]
\centering
\footnotesize
\begin{tabularx}{0.88\linewidth}{@{}p{2.5cm} X@{}}
\toprule
\textbf{Symbol} & \textbf{Meaning} \\
\midrule
$o,y_0,f,B$ & Campaign objective, initial artifact, evaluator, and resource budget \\
$\tau_t$ & Mission trajectory $(s_t,a_t,y_t,e_t,r_t)$ \\
$\rho_I(T)$ & Dense-intelligence density over wall-clock horizon $T$ \\
$\dot N_{\mathrm{tok}}$ & Rate of token-mediated reasoning and action \\
$\eta_r,\eta_a,\eta_v$ & Reasoning, action, and verification conversion factors \\
$H_t$ & Runtime state $\{M_t,S_t,A_t,V_t,R_t,Q_t\}$ \\
$E_t$ & Results admitted from mission $t$ \\
$\theta_t$ & Underlying model parameters, fixed in this report \\
$C_t$ & Effective retained capability set \\
$W_t,D_t(d)$ & Capability breadth and depth in domain $d$ \\
$\epsilon$ & Acceptance threshold in the idealized capability-admission rule \\
$V_R(T)$ & Conceptual accumulated research value over horizon $T$ \\
$\Delta I,p_{\mathrm{valid}},\eta_{\mathrm{reuse}}$ & Information gain, validity probability, and reuse efficiency \\
$D_{\mathrm{process}}$ & Retained states, actions, measurements, feedback, and runtime-state deltas \\
\bottomrule
\end{tabularx}
\caption{Notation used in the main text.}
\end{table}

The paper uses three figures, each with a distinct evidentiary purpose. The first-page
teaser, \code{argus\_teaser}, combines the exact objective--runtime--review--state
flow with protocol-scoped public result cards. It is an HTML/CSS vector composition
generated from the committed result JSON by
\srcpath{figures/build\_argus\_teaser.py}; its figure brief, source hashes, editable
HTML, and renderer are recorded in \srcpath{figures/argus\_teaser.json}. The two
empirical figures, \code{public\_results} and \code{paper\_portfolio}, are generated
by \srcpath{figures/build\_report\_figures.py} and digested in
\srcpath{figures/REPORT\_FIGURES.json}. The remaining structural assets in the
repository are intentionally not included: their content is represented more
precisely by equations, tables, and interface definitions in this version.
